\begin{textsample}{NEG Dim 2 – 2022\_c4 – Score 18.00 – 2022\_c4/enade\_2022\_ci\_cont\_9597f67a94.txt}  \label{ex:f2_neg_406}
CAPÍTULO 5 PERCEPÇÃO DA PROVA Nas análises feitas neste capítulo, constam as percepções dos concluintes da área de Ciências Contábeis sobre a prova aplicada no Enade/2022.
Essas percepções foram mensuradas por meio de nove questões que avaliaram desde o grau de dificuldade da prova até o tempo gasto para concluí -la.
As percepções sobre a prova foram relacionadas ao desempenho dos estudantes, à grande região de funcionamento do curso, à categoria administrativa e ao tipo de organização acadêmica da IES.
O questionário de percepção da prova encontra -se ao final do Anexo VIII, no qual está a reprodução do exame.
Na apresentação dos dados relativos às nove questões sobre as percepções a respeito da prova, o desempenho dos estudantes foi classificado em quatro quartos.
Para tanto, esse desempenho foi ordenado de forma ascendente.
O percentil 25, P25, também conhecido como primeiro quartil, é a nota de desempenho que deixa um quarto ( 25 \% ) dos valores observados abaixo e três quartos acima.
O quarto inferior de desempenho é composto pelas notas abaixo do primeiro quartil.
Já o percentil 75, P75, também conhecido como terceiro quartil, é o valor para o qual há três quartos ( 75 \% ) dos dados abaixo e um quarto acima dele.
O quarto superior de desempenho é composto pelas notas iguais ou acima do terceiro quartil.
O percentil 50, P50, também conhecido como mediana, é o valor que divide as notas em dois conjuntos de igual tamanho.
O segundo quarto inclui os valores entre o primeiro quartil ( P25 ) e a mediana.
O terceiro quarto contém os valores entre a mediana ( P50 ) e o terceiro quartil ( P75 ).
Vale ressaltar que percentis, quartis e medianas são pontos que, não obrigatoriamente, pertencem ao conjunto original de dados, ao passo que os quartos são subconjuntos dos dados originais.
Na Figura 5.1, apresenta -se uma ilustração dos quatro quartos descritos.
A seguir, serão apresentados gráficos com resultados selecionados, relativos às nove questões avaliadas por grupos de estudantes.
Nas barras dos gráficos, apresenta -se o percentual de estudantes que assinalaram uma das opções ou a soma das porcentagens daqueles que assinalaram duas ( ou três ) delas.
Por exemplo, para as questões 1 e 2, nos gráficos, é apresentada a porcentagem total de participantes que assinalaram as opções ( D ) “ Difícil ” e ( E ) “ muito Difícil ”.
Em cada barra, foram assinalados, também, os extremos do intervalo de confiança de 95 \% como linhas verticais unidas por uma linha horizontal na forma da letra H maiúscula.
O estimador de um parâmetro com um certo nível de confiança ( por exemplo, 95,0 \% ) deve conter o parâmetro no intervalo de confiança em 95 \% das vezes.
Na comparação entre os estimadores dos parâmetros de duas classes de uma dada categoria ( por exemplo : Norte e Nordeste nas grandes regiões, de primeiro e último quarto dentro de desempenho, de IES públicas e privadas, ou de universidades e faculdades ), associados aos seus respectivos intervalos de confiança, diz -se que não há diferença estatisticamente significativa entre os parâmetros das duas categorias quando há interseção entre os intervalos de confiança, e que há diferença, se os intervalos de confiança são disjuntos ( para maiores informações vide Glossário ).
Nas Tabelas do Anexo V, são apresentados os valores absolutos e a distribuição percentual9 das alternativas válidas das nove questões, segundo o mesmo recorte de desempenho dos estudantes e da grande região de funcionamento do curso, de categoria administrativa e do tipo de organização acadêmica da IES. 5.1 GRAU DE DIFICULDADE DA PROVA Nesta seção são mostradas as percepções dos estudantes quanto ao grau de dificuldade da prova Enade/2022, e que foram analisadas considerando sua vinculação ao Componente de Formação Geral e ao Componente Específico. 5.1.1 COMPONENTE DE FORMAÇÃO GERAL O Gráfico 5.1 apresenta as respostas obtidas para a Questão 1 : “ Qual o grau de dificuldade desta prova na parte de formação geral?
De acordo com os dados do Gráfico 5.1, a região Sul apresentou o maior percentual de estudantes que consideraram a prova “ Difícil ” ou “ Muito Difícil ”, com 40,6 \%, seguida da região Centro-Oeste com 38,8 \%.
O menor percentual foi apresentado na região Sudeste, 32,8 \%.
Na região Nordeste, a média foi de 33,5 \%, na Norte, 37,3 \%.
O Gráfico 5.2 traz mais resultados sobre a questão tratada no gráfico anterior, porém por quartos de desempenho.
Os dados do Gráfico 5.2 mostram que o percentual de estudantes que consideraram o componente de Formação Geral da prova “ Difícil ” ou “ Muito Difícil ” foi decrescente em função dos quartos de desempenho : 41,2 \% no 1º quarto ; e 29,4 \%, no 4º quarto – o grupo de melhor desempenho na prova.
Nos quartos de desempenho intermediários, a proporção de estudantes que consideraram a prova “ Díficil ” ou “ Muito Díficil ” foi de 38,1 \% no 2º quarto ; e de 35,3 \% no 3º quarto.
As diferenças entre todos os quartos de desempenho são estatisticamente significativas.
O Gráfico 5.3 traz outros resultados também em relação à mesma questão retratada no gráfico anterior, porém por organização acadêmica das IES.
Como mostram os dados do Gráfico 5.3, o percentual de estudantes que consideraram o componente de Formação Geral da prova “ Difícil ” ou “ Muito Difícil ” foi maior entre os estudantes de CEFET/IF, alcançando 57,8 \%, enquanto, entre os estudantes das universidades, a organização acadêmica com a menor incidência, o percentual alcançou 34,6 \%.
O Gráfico 5.4 apresenta os dados por categoria administrativa das IES.
Os dados do Gráfico 5.4 mostram que o percentual de estudantes que consideraram o componente de Formação Geral da prova “ Difícil ” ou “ Muito Difícil ” foi maior entre os estudantes de cursos de IES privadas ( 37,4 \% ), com diferença estatisticamente significativa da proporção registrada em IES públicas ( 28,5 \% ). 5.1.2 COMPONENTE DE CONHECIMENTO ESPECÍFICO O Gráfico 5.5 apresenta as respostas obtidas para a Questão 2 : “ Qual o grau de dificuldade desta prova na parte de Componente Específico? ”.
De acordo com os dados do Gráfico 5.15, a região Sul apresentou o maior percentual de estudantes que consideraram o componente específico da prova “ Difícil ” ou “ Muito Difícil ”, com 56,3 \%, seguida da região Norte com 53,9 \%.
Na região Nordeste, a média foi de 52,7 \% ; na Sudeste, 51,4 \% ; e na região Centro-Oeste, 53,4 \%.
O Gráfico 5.6 traz mais resultados sobre a questão tratada no gráfico anterior, porém por quartos de desempenho.
Os dados do Gráfico 5.6 mostram que o percentual de estudantes que consideraram o componente Específico da prova “ Difícil ” ou “ Muito Difícil ” foi crescente em função dos quartos de desempenho : 49,1 \%, no 1.º quarto e 55,8 \%, no 4.º quarto, o grupo de melhor desempenho na prova.
Nos quartos de desempenho intermediários, a proporção de estudantes que consideraram a prova “ Difícil ” ou “ Muito Difícil ” foi de 53,1 \%, no 2.º quarto, e de 54,9 \% no 3.º quarto.
O Gráfico 5.7 traz outros resultados também em relação à mesma questão retratada no gráfico anterior, porém por organização acadêmica das IES.
Como mostram os dados do Gráfico 5.7, o percentual de estudantes que consideraram o componente de formação geral da prova “ Difícil ” ou “ Muito Difícil ” foi maior entre os estudantes de CEFET/IF, alcançando 68,9 \%, enquanto entre os estudantes das faculdades, a organização acadêmica com a menor incidência, o percentual alcançou 52,4 \%.
No mesmo gráfico, é possível observar que a diferença de CEFET/IF para as demais organizações acadêmicas é estatisticamente significativa.
O Gráfico 5.8 apresenta os dados por categoria administrativa das IES.
Os dados do Gráfico 5.8 mostram que o percentual de estudantes que consideraram o componente Específico da prova “ Difícil ” ou “ Muito Difícil ” foi maior entre os estudantes de cursos de IES privadas ( 53,5 \% ), enquanto a proporção registrada para as IES públicas foi de 51,8 \%. 5.2 EXTENSÃO DA PROVA EM RELAÇÃO AO TEMPO TOTAL O Gráfico 5.9 apresenta as respostas obtidas para a Questão 3 : “ Considerando a extensão da prova, em relação ao tempo total, você considera que a prova foi longa ou muito longa? ”.
Em relação aos dados do Gráfico 5.9, entre as grandes regiões, os que avaliaram a prova como longa ou muito longa, em relação ao tempo total destinado à sua resolução, atingiram uma porcentagem que variou de 32,6 \% na região Norte, até 42,4 \%, na região Sul.
O Gráfico 5.10 traz mais resultados sobre a questão tratada no gráfico anterior, porém por quartos de desempenho.
No Gráfico 5.10, observa -se que a proporção de estudantes que consideraram a prova longa ou muito longa foi crescente em função dos quartos de desempenho : 40 \% no 1.º quarto ; e 41,3 \% no 4.º quarto – o grupo de melhor desempenho na prova.
Nos quartos de desempenho intermediários, a proporção de estudantes que consideraram a prova longa ou muito longa foi de 38,9 \% no 2.º quarto ; e de 39,3 \% no 3.º quarto.
O Gráfico 5.11 traz outros resultados também em relação à mesma questão retratada no gráfico anterior, porém por organização acadêmica das IES.
De acordo com os dados do Gráfico 5.11, considerando -se a extensão da prova, o percentual de estudantes que avaliaram a prova como longa ou muito longa, em relação ao tempo total destinado à sua resolução, foi maior entre os estudantes das faculdades, alcançando 41,6 \%, enquanto, entre os estudantes de universidades, a organização acadêmica com a menor incidência, esse índice atingiu 39,6 \%.
Entre os participantes de centros universitários, o índice atingiu 39,9 \%, e entre os de CEFET/IF, 40,0 \%.
O Gráfico 5.12 apresenta os dados por categoria administrativa das IES.
Os dados do Gráfico 5.12 mostram que o percentual de estudantes que consideraram a prova longa ou muito longa, em relação ao tempo total destinado à sua resolução, foi maior entre os estudantes de cursos de IES privadas ( 40 \% ) do que entre os estudantes de IES públicas ( 39,3 \% ), não havendo diferença estatisticamente significativa da proporção entre as duas categorias administrativas das IES. 5.3 COMPREENSÃO DOS ENUNCIADOS DAS QUESTÕES Também organizados em relação ao Componente de Formação Geral e ao Componente Específico, os dados apresentados nesta seção dizem respeito à forma como os estudantes compreenderam os enunciados das questões da prova, considerando as formulações expressas nesses mesmos enunciados. 5.3.1 COMPONENTE DE FORMAÇÃO GERAL O Gráfico 5.13 apresenta as respostas obtidas para a Questão 4 : “ Os enunciados das questões da prova na parte de Formação Geral estavam claros e objetivos? ”.
Os resultados apresentados no Gráfico 5.13 mostram que, na análise regional, a porcentagem de estudantes que avaliaram que todos ou a maioria dos enunciados das questões do Componente de Formação Geral estavam “ claros e objetivos ” variou de 75,6 \% na região Norte a 76,8 \% na região Sudeste.
Nas demais regiões, as porcentagens foram de : 76 \% para a região Nordeste, 75,7 \% para a região Sul e 76,6 \% para a região Centro-Oeste.
O Gráfico 5.14 traz mais resultados sobre a questão tratada no gráfico anterior, porém por quartos de desempenho.
De acordo com os dados do Gráfico 5.14, segundo o desempenho, observa -se que a proporção dos que emitiram essa opinião cresce, conforme o desempenho aumenta, com diferenças estatisticamente significativas entre todos os quartos de desempenho.
No quarto superior, a clareza e a objetividade de todos ou da maioria dos enunciados das questões foram percebidas por 83,1 \%.
O Gráfico 5.15 traz outros resultados também em relação à mesma questão retratada no gráfico anterior, porém por organização acadêmica das IES.
Os dados do Gráfico 5.15 mostram que a proporção dos presentes que consideraram que todos ou a maioria dos enunciados das questões do Componente de Formação Geral estavam “ claros e objetivos ” foi maior entre os estudantes de CEFET/IF, chegando a 82,2 \%, enquanto entre os de faculdades, a organização acadêmica com a menor incidência, esse índice atingiu 75,6 \%.
Entre os estudantes de centros universitários, o índice atingiu 76,0 \%, e entre os de universidades, 76,6 \%.
O Gráfico 5.16 apresenta os dados por categoria administrativa das IES.
Os dados do Gráfico 5.16 mostram que o percentual de estudantes que consideraram que todos ou a maioria dos enunciados das questões do Componente de Formação Geral estavam “ claros e objetivos ” foi maior entre os estudantes de cursos de IES públicas ( 80,8 \% ) quando comparado aos estudantes de IES privadas ( 75,5 \% ), e com diferença estatisticamente significativa da proporção entre as categorias administrativas. 5.3.2 COMPONENTE ESPECÍFICO O Gráfico 5.17 apresenta as respostas obtidas para a Questão 5 : “ Os enunciados das questões da prova na parte de Componente Específico estavam claros e objetivos? ”.
A maioria dos estudantes de todas as grandes regiões brasileiras considerou “ claros e objetivos ” todos ou a maioria dos enunciados das questões do Componente Específico da prova, percentual sempre maior ou igual a 70,1 \%, para a região Norte, chegando ao maior percentual de 72,5 \% para a região Sul.
O Gráfico 5.18 traz mais resultados sobre a questão tratada no gráfico anterior, porém por quartos de desempenho.
Os dados do Gráfico 5.18 mostram que a proporção de estudantes que consideraram os enunciados das questões “ claros e objetivos ” apresenta uma tendência crescente em relação ao aumento de desempenho : mais elevada no quarto superior ( 78,3 \% ), se comparada ao quarto inferior de desempenho ( 66,0 \% ).
As diferenças entre todos os quartos de desempenho são estatisticamente significativas.
O Gráfico 5.19 traz outros resultados também em relação à mesma questão retratada no gráfico anterior, porém por organização acadêmica das IES.
Os dados do Gráfico 5.19 mostram que a proporção dos presentes que consideraram que todos ou a maioria dos enunciados das questões do Componente Específico estavam “ claros e objetivos ” foi maior entre os estudantes de CEFET/IF, atingindo 77,8 \%.
Entre os estudantes de faculdades, organização acadêmica com a menor incidência, essa proporção atingiu 70,6 \%, sendo intermediárias as proporções alcançadas entre os estudantes de universidades ( 72,1 \% ) e de centros universitários ( 71,1 \% ).
No Gráfico 5.19, é possível observar que as diferenças de universidades para faculdades e centros universitários são estatisticamente significativas.
O Gráfico 5.20 apresenta os dados por categoria administrativa das IES.
Os dados do Gráfico 5.20 mostram que o percentual de estudantes que consideraram que todos ou a maioria dos enunciados das questões do Componente Específico estavam “ claros e objetivos ” foi maior entre os estudantes de cursos de IES públicas ( 75,4 \% ) que entre os de IES privadas ( 70,9 \% ), com diferença estatisticamente significativa da proporção registrada entre elas. 5.4 SUFICIÊNCIA DAS INFORMAÇÕES/INSTRUÇÕES FORNECIDAS O Gráfico 5.21 apresenta as respostas obtidas para a Questão 6 : “ As informações/instruções fornecidas para a resolução das questões foram suficientes para resolvê -las? ” De acordo com os dados do Gráfico 5.21 e quanto à distribuição de respondentes pelas grandes regiões, observa -se que a proporção de estudantes que consideraram as informações/instruções fornecidas até excessivas ou suficientes em todas ou na maioria das questões foi sempre igual ou superior ou igual a 76,9 \%, chegando a 79,8 \%, na região Sudeste.
O Gráfico 5.22 traz resultados sobre a questão tratada no gráfico anterior, porém por quartos de desempenho.
De acordo com os dados do Gráfico 5.22, levando -se em conta o desempenho dos participantes, observa -se uma tendência crescente das proporções de participantes que avaliaram as informações/instruções como até excessivas ou suficientes em todas ou na maioria das questões, havendo diferenças estatisticamente significativas entre todos os quartos de desempenho.
O percentual foi mais elevado no quarto superior ( 84,0 \% ).
Já no quarto inferior, a suficiência das informações/instruções declarada como até excessiva, em todas ou na maioria das questões foi percebida por 74,3 \% dos respondentes.
O Gráfico 5.23 traz outros resultados também em relação à mesma questão retratada no gráfico anterior, porém por organização acadêmica das IES.
De acordo com os dados do Gráfico 5.23, considerando -se as informações/instruções fornecidas para a resolução das questões ( Questão 6 ), a proporção que considerou que essas eram até excessivas ou suficientes em todas ou na maioria das questões foi maior entre os estudantes das universidades, totalizando 79 \%.
Os CEFET/IF foram a organização acadêmica em que houve a menor proporção : 75,6 \%.
O Gráfico 5.24 apresenta os dados por categoria administrativa das IES.
De acordo com os dados do Gráfico 5.24, o percentual de estudantes que consideraram que as informações/instruções fornecidas para a resolução das questões eram até excessivas ou suficientes em todas ou na maioria das questões foi maior entre os estudantes de cursos de IES públicas ( 83,4 \% ), diferença estatisticamente significativa da proporção registrada entre os de IES privadas ( 77,6 \% ). 5.5 DIFICULDADE ENCONTRADA AO RESPONDER À PROVA O Gráfico 5.25 apresenta as respostas obtidas para a Questão 7 : “ Você se deparou com alguma dificuldade ao responder à prova?
Qual? ” Nos Gráficos 5.25 a 5.28, apresentados de forma sequenciada mais adiante, são expostos os percentuais de estudantes que apontaram o “ Desconhecimento do conteúdo ” como “ dificuldade ao responder à prova ”.
O Gráfico 5.25 apresenta as respostas obtidas sobre o “ Desconhecimento do conteúdo ” como a principal “ dificuldade ao responder à prova ”.
De acordo com os dados o Gráfico 5.25, o percentual de inscritos e presentes que apontaram o desconhecimento do conteúdo como “ dificuldade ao responder à prova ” não superou 20,0 \%.
Os percentuais variaram de 17,9 \% na região Nordeste a 20 \% na região Norte.
O Gráfico 5.26 traz resultados sobre o ponto tratado no gráfico anterior, porém por quartos de desempenho.
De acordo com os dados do Gráfico 5.26, em relação aos quartos de desempenho, o desconhecimento do conteúdo foi a opção escolhida por 17,3 \% dos estudantes do quarto inferior e por 20,0 \% do quarto superior, caracterizando uma tendência crescente.
O Gráfico 5.27 traz outros resultados também em relação à mesma questão retratada no gráfico anterior, porém por organização acadêmica das IES.
De acordo com os dados do Gráfico 5.27, na análise por organização acadêmica, o percentual de inscritos e presentes que apontaram o Desconhecimento do conteúdo como “ dificuldade ao responder à prova ” não superou 22,2 \%.
Os percentuais variaram de 17,5 \%, nas faculdades, a 22,2 \% nos CEFET/IF.
O Gráfico 5.28 apresenta os dados por categoria administrativa das IES.
De acordo com os dados do Gráfico 5.28, em relação às categorias administrativas, o desconhecimento do conteúdo foi a opção escolhida por 20,1 \% dos estudantes, nas IES públicas, e por 18,1 \%, nas IES privadas.
São observadas diferenças estatisticamente significativas entre essas duas categorias. 5.6 CONTEÚDOS DAS QUESTÕES OBJETIVAS DA PROVA O Gráfico 5.29 apresenta as respostas obtidas para a Questão 8 : “ Considerando apenas as questões objetivas da prova, você percebeu que não estudou ainda a maioria desses conteúdos? ”.
Na análise por grande região, a proporção de respondentes que escolheram a opção não estudou ainda a maioria desses conteúdos foi pequena.
Observa -se que, na região Norte ( 7,8 \% ), a proporção foi maior.
Na região Sudeste, a situação foi inversa, com proporção mais baixa, correspondendo a 5,5 \%.
Na sequência, têm -se as regiões Sul, Nordeste e Centro-Oeste com proporções de 6,0 \%, 6,6 \% e 6,9 \%, respectivamente.
O Gráfico 5.30 traz resultados sobre a questão tratada no gráfico anterior, porém por quartos de desempenho.
De acordo com os dados do Gráfico 5.30, considerando -se separadamente as opiniões de estudantes dos quatro quartos de desempenho, observa -se que, no quarto inferior, 9,6 \% ofereceram como resposta que não estudaram ainda a maioria desses conteúdos, sendo 3,6 \% os do quarto superior com a mesma resposta.
As proporções são decrescentes com o desempenho.
A diferença entre os estudantes que optaram por esse motivo de dificuldade entre todos os quartos de desempenho é estatisticamente significativa.
O Gráfico 5.31 traz outros resultados também em relação à mesma questão retratada no gráfico anterior, porém por organização acadêmica das IES.
De acordo com os dados do gráfico 5.31, na análise por organização acadêmica, a proporção de respondentes que escolheram a opção não estudou ainda a maioria desses conteúdos foi pequena.
Observa -se que as proporções nos CEFET/IF ( 2,2 \% ) e nas faculdades ( 6 \% ) foram as menores.
Nas universidades e nos centros universitários, a proporção foi mais alta, 6,2 \% e 6,4 \%, respectivamente.
O Gráfico 5.32 apresenta os dados por categoria administrativa das IES.
De acordo com os dados do Gráfico 5.32, em relação às categorias administrativas, observa -se que, nas IES públicas, 4,2 \% ofereceram como resposta que não estudou ainda a maioria desses conteúdos, sendo 6,6 \% nas IES privadas com a mesma resposta.
A diferença entre os estudantes que optaram por esse motivo de dificuldade nas duas categorias é estatisticamente significativa. 5.7 TEMPO GASTO PARA CONCLUIR A PROVA O Gráfico 5.33 apresenta as respostas obtidas para a Questão 9 : “ Qual foi o tempo gasto por você para concluir a prova? ” Considerando -se as cinco grandes regiões brasileiras, os percentuais dos que utilizaram entre duas e quatro horas para finalizar a prova, na região Norte, foi de 82,4 \%.
Nas regiões Sul, Sudeste e Centro-Oeste, os percentuais de estudantes que dispensaram entre duas e quatro horas para concluir a prova ficaram próximos : 69,5 \%, 69,3 \% e 68,2 \%, respectivamente.
A região Nordeste ficou com o percentual de 74,1 \%.
O Gráfico 5.34 traz resultados sobre a questão tratada no gráfico anterior, porém por quartos de desempenho.
De acordo com os dados do Gráfico 5.34, uma vez analisadas as alternativas escolhidas pelos estudantes que se situam nos diferentes quartos de desempenho, observam -se uma tendência crescente : uma proporção maior de participantes no quarto superior ( 81,9 \% ) declarou ter gasto entre duas e quatro horas para concluir a prova quando comparada com a proporção de participantes nos quartos inferiores.
Há diferenças estatisticamente significativas entre todos os quartos de desempenho.
O Gráfico 5.35 traz outros resultados também em relação à mesma questão retratada no gráfico anterior, porém por organização acadêmica das IES.
De acordo com os dados do Gráfico 5.35, na análise por organização acadêmica, os percentuais dos que utilizaram Entre duas e quatro horas para finalizar a prova foram de 72,4 \% nos centros universitários, 77,5 \% nas faculdades e 88,9 \% nos CEFET/IF.
O menor percentual foi entre os estudantes das universidades, com 68,3 \% dos alunos necessitando de duas a quatro horas para concluir a prova.
Como mostrado no Gráfico 5.35, há diferença estatisticamente significativa entre as proporções de todos os tipos de organização acadêmica.
O Gráfico 5.36 apresenta os dados por categoria administrativa das IES.
De acordo com os dados do Gráfico 5.36, em relação às categorias administrativas, observa -se que, nas IES públicas, 72,8 \% declararam ter gasto Entre duas e quatro horas para concluir a prova, sendo 70,3 \% nas IES privadas com a mesma resposta.
A diferença entre os estudantes que optaram por esse motivo de dificuldade nas duas categorias é estatisticamente significativa.
Os resultados do Enade 2022 apresentados neste capítulo contribuem para o entendimento de que a consideração da percepção dos estudantes concluintes sobre a prova do Enade é importante tanto para a compreensão do seu desempenho quanto para a reflexão sistemática e constante a respeito de determinadas realidades relacionadas a esse mesmo desempenho, no país.
Por outro lado, como todos os elementos do trabalho pedagógico institucional mostram -se diretamente relacionados ao desempenho dos estudantes, também ganha relevância conhecer e analisar os resultados da percepção dos coordenadores de curso sobre a mesma prova.
As informações expressas por estudantes e coordenadores sobre questões diversas associadas aos dois componentes da prova – Formação Geral e Conhecimento Específico – podem agregar valor em ações institucionais voltadas ao aperfeiçoamento do instrumento avaliativo, considerando os vários elementos envolvidos em sua elaboração.
Neste contexto, é relevante que os resultados das percepções de estudantes e coordenadores sobre a prova do Enade ganhem peso como um elemento a mais nas análises da oferta dos cursos nas modalidades presencial e a distância, em nível nacional.
Espera -se, assim, que, em uma visão formativa de avaliação, o conhecimento e a análise desses resultados possam vir a contribuir para a construção de alternativas institucionais que venham a induzir a qualidade da educação superior.

% matched lemmas: 
\end{textsample}
